%% kapa: exam=1 round=37 year=2026 period=2 form=A subject=accounting subject_ko=회계학 seq=17 number=57
%% source: Q-Net 감정평가사 기출문제 (cst00309) artlSeq=5258924 file=raw/1008/5258924/2026년 제37회 감정평가사 1차 2교시.pdf
%% conversion: kordoc 4.24.51 -> tools/gen1.py (auto); review=proofread-sonnet; flags=none
\begin{question}{57}
주당이익 회계처리에 관한 설명으로 옳은 것은?
\begin{choices}
  \item 할증배당우선주의 당초 할인발행차금이나 할증발행차금은 유효이자율법으로 상각하여 이익잉여금에 가감하고 주당이익을 계산할 때 우선주 배당금으로 처리한다.
  \item 희석주당이익 계산시 증분주식 1주당 이익이 가장 큰 희석성 잠재적보통주를 증분주식 1주당 이익이 상대적으로 작은 희석성 잠재적보통주보다 먼저 포함시킨다.
  \item 재매입가능성이 높은 조건부발행보통주는 가중평균하여 기본주당이익을 계산하기 위한 보통주식수에 포함한다.
  \item 자본금전입, 무상증자, 주식분할의 경우에는 추가로 대가를 받지 않고 기존 주주에게 보통주를 발행하므로 유통보통주식수는 증가하지 않는다.
  \item 당기에 발행된 희석성 잠재적보통주는 회계기간의 기초에 전환된 것으로 본다.
\end{choices}
\end{question}
\answer{1}
